% chapters/04c-case-study.tex — Case study: บ้านพักอาศัย 100 ตรม. BOQ
\section{กรณีศึกษา: การประยุกต์ใช้ระบบกับโครงการก่อสร้างจริง}
\label{sec:case-study}

เพื่อแสดงประโยชน์ของระบบในงานวิศวกรรมโยธาในเชิงปฏิบัติ ผู้วิจัยได้นำระบบ
ไปทดสอบกับโครงการก่อสร้างบ้านพักอาศัย 1 ชั้น พื้นที่ใช้สอย 100 ตร.ม.
ในจังหวัดนครราชสีมา (ภาคตะวันออกเฉียงเหนือ) โดยเปรียบเทียบกระบวนการ
ประมาณราคาของวิศวกรประมาณราคาด้วยวิธีดั้งเดิม (manual estimation)
กับการใช้ระบบที่พัฒนาขึ้นในงานวิจัยนี้.

\subsection{ขอบเขตของกรณีศึกษา}

\begin{itemize}[leftmargin=2em]
  \item \textbf{ประเภทโครงการ}: บ้านพักอาศัย 1 ชั้น โครงสร้างคอนกรีตเสริมเหล็ก
        ผนังก่ออิฐมอญฉาบปูน หลังคา metal sheet.
  \item \textbf{พื้นที่}: 100 ตร.ม.\ (รวมพื้นที่ส่วนรับรอง 60 ตร.ม.\
        ห้องนอน 30 ตร.ม.\ และห้องน้ำ 10 ตร.ม.).
  \item \textbf{ที่ตั้ง}: อำเภอเมือง จังหวัดนครราชสีมา.
  \item \textbf{ขอบเขตการประมาณ}: เฉพาะวัสดุก่อสร้างหลัก 5 หมวด
        ตามประเภทงานที่ระบบรองรับ — งานก่ออิฐ งานเสาเอ็น/คานทับหลัง
        งานเหล็กเสริม งานทาสี และงานบุกระเบื้อง.
\end{itemize}

\subsection{ปริมาณงานที่ใช้ในการคำนวณ}

ปริมาณงาน (work quantity) ที่ป้อนเข้าระบบประเมินจากแบบสถาปัตยกรรม
และแบบโครงสร้างของบ้านตัวอย่าง ดังตารางที่ \ref{tab:case-quantity}.

\begin{table}[H]
  \centering
  \footnotesize
  \caption{ปริมาณงานที่ใช้ในการคำนวณ BOQ ของบ้าน 100 ตร.ม.}
  \label{tab:case-quantity}
  \begin{tabular}{llrl}
    \toprule
    หมวดงาน & รายการ & ปริมาณ & หน่วย \\
    \midrule
    งานก่ออิฐ        & ผนังภายใน-ภายนอก                  & 240   & ตร.ม. \\
    งานเสาเอ็น       & เสาเอ็น/คานทับหลัง $0.10\times0.15$ & 86    & ม.    \\
    งานเหล็กเสริม    & DB12 (เสา-คาน)                    & 1{,}450 & กก.   \\
    งานเหล็กเสริม    & RB6 (ปลอก stirrup)                & 380   & กก.   \\
    งานทาสี          & สีรองพื้น $+$ สีจริง 2 เที่ยว       & 480   & ตร.ม. \\
    งานบุกระเบื้อง   & ห้องน้ำ-ห้องครัว                    & 32    & ตร.ม. \\
    \bottomrule
  \end{tabular}
\end{table}

\subsection{ผลการเปรียบเทียบ}

ตารางที่ \ref{tab:case-compare} แสดงผลการประมาณต้นทุนวัสดุของแต่ละหมวดงาน
จากสองวิธี — manual estimation โดยอ้างอิงเฉพาะ CGD baseline เทียบกับ
การใช้ระบบที่อ้างอิงราคา Modern Trade เฉลี่ย 8 แหล่ง ในภาคตะวันออกเฉียงเหนือ.

\begin{table}[H]
  \centering
  \footnotesize
  \caption{เปรียบเทียบต้นทุนวัสดุของบ้าน 100 ตร.ม.\ ระหว่างวิธีดั้งเดิมกับการใช้ระบบ}
  \label{tab:case-compare}
  \begin{tabular}{lrrr}
    \toprule
    หมวดงาน & Manual (CGD only) & ระบบ (MT, อีสาน) & ส่วนต่าง \\
            & (บาท) & (บาท) & \\
    \midrule
    งานก่ออิฐ            &  72{,}000 &  85{,}500 & $+$18.8\% \\
    งานเสาเอ็น/คาน        &  18{,}920 &  22{,}620 & $+$19.6\% \\
    งานเหล็กเสริม         &  46{,}980 &  56{,}890 & $+$21.1\% \\
    งานทาสี              &  21{,}600 &  26{,}160 & $+$21.1\% \\
    งานบุกระเบื้อง        &  18{,}560 &  23{,}040 & $+$24.1\% \\
    \midrule
    \textbf{รวมต้นทุนวัสดุ} & \textbf{178{,}060} & \textbf{214{,}210} & \textbf{$+$20.3\%} \\
    \bottomrule
  \end{tabular}
\end{table}

ผลลัพธ์ชี้ให้เห็นว่าการประมาณราคาด้วย CGD baseline เพียงอย่างเดียว
ทำให้ต้นทุนวัสดุของโครงการต่ำกว่าราคาตลาดจริงประมาณ 20.3\%
($\Delta = 36{,}150$ บาท) ส่วนต่างนี้สอดคล้องกับผลการเปรียบเทียบเชิง
สถิติในหัวข้อ \ref{sec:validation} และเป็นความเสี่ยงที่ผู้รับเหมา
ต้องแบกรับหากเข้าประกวดราคาโดยอ้างอิงเฉพาะราคาภาครัฐ.

\subsection{การวัดเวลาในการจัดทำ BOQ}

นอกจากความถูกต้องของราคา ผู้วิจัยได้จับเวลาในการจัดทำ BOQ ของกรณีศึกษานี้
ด้วยสองวิธี ทำซ้ำ 5 ครั้ง โดยผู้ใช้รายเดียวกันที่มีประสบการณ์การประมาณ
ราคาในระดับเดียวกัน. ผลแสดงในตารางที่ \ref{tab:case-time}.

\begin{table}[H]
  \centering
  \footnotesize
  \caption{เวลาเฉลี่ยในการจัดทำ BOQ บ้าน 100 ตร.ม.\ ($n = 5$)}
  \label{tab:case-time}
  \begin{tabular}{lcc}
    \toprule
    วิธี & เวลาเฉลี่ย (นาที) & ส่วนเบี่ยงเบนมาตรฐาน \\
    \midrule
    Manual (สืบค้น 8 เว็บแยก $+$ Excel)    & 47.6 & 5.8 \\
    ระบบ (ป้อนปริมาณ $\to$ export)        &  8.4 & 1.2 \\
    \midrule
    \textbf{เวลาที่ลดลง}                  & \textbf{$-$82.4\%} &  \\
    \bottomrule
  \end{tabular}
\end{table}

ระบบช่วยลดเวลาในการจัดทำ BOQ ลง 82.4\% (จาก 47.6 นาที เหลือ 8.4 นาที
ต่อโครงการ) ในขณะเดียวกันยังให้ราคาที่สะท้อนตลาดจริงในภาคตะวันออกเฉียงเหนือ
ซึ่งเป็นปัจจัยที่ราคาภาครัฐเพียงอย่างเดียวไม่ครอบคลุม. กรณีศึกษานี้แสดง
ให้เห็นบทบาทของระบบในฐานะเครื่องมือสนับสนุนวิศวกรประมาณราคาทั้งใน
มิติของความถูกต้อง (accuracy) และมิติของประสิทธิภาพการทำงาน
(productivity).
